% --- Archon physics formalization source begin ---
% archon:physics
% archon:covers PhyXMiniProblems/problem_phyx_mini_0235.lean
% archon:source-report reports/phyx_mini/problem_phyx_mini_0235.source.json

\chapter{Physics problem phyx\_mini\_0235}
\label{ch:PhyXMiniProblems_problem_phyx_mini_0235}

\paragraph{Problem source.}
\#\# Physical scenario

One end of a light spring with force constant \textbackslash{}( k = 100\textbackslash{},\textbackslash{}mathrm\{N/m\} \textbackslash{}) is attached to a vertical wall. A light string is tied to the other end of the horizontal spring. As shown in figure, the string changes from horizontal to vertical as it passes over a pulley of mass \textbackslash{}( M \textbackslash{}) in the shape of a solid disk of radius \textbackslash{}( R = 2.00\textbackslash{},\textbackslash{}mathrm\{cm\} \textbackslash{}). The pulley is free to turn on a fixed, smooth axle. The vertical section of the string supports an object of mass \textbackslash{}( m = 200\textbackslash{},\textbackslash{}mathrm\{g\} \textbackslash{}). The string does not slip at its contact with the pulley. The object is pulled downward a small distance and released.The pulley radius is doubled to \textbackslash{}( R = 4.00\textbackslash{},\textbackslash{}mathrm\{cm\} \textbackslash{}).

\#\# Question

What is the highest possible value of the angular frequency of oscillation of the object?

\#\# Answer choices

- A: A: \$20.4\textbackslash{}

- B: B: \$23.4\textbackslash{}

- C: C: \$21.4\textbackslash{}

- D: D: \$22.4\textbackslash{}

\#\# Figure caption (auxiliary; use the image as primary evidence)

The image illustrates a mechanical setup involving a spring, pulley, and blocks. Here's the detailed description:

- **Spring (k):** On the left side, there is a spring attached to a vertical wall. The spring constant is labeled \textbackslash{}( k \textbackslash{}).

- **Pulley System:** The spring is connected to a rope that goes over a pulley. The pulley is mounted on a horizontal surface and is depicted with a blue triangular support.

- **Rope:** The rope runs over the pulley and extends downward, holding a mass labeled \textbackslash{}( m \textbackslash{}). 

- **Block M:** The point on the rope that's connected to the spring is labeled \textbackslash{}( M \textbackslash{}), denoting another mass or a point of interest.

- **Distances:** There are two horizontal distances labeled on the setup; one is the extension of the spring \textbackslash{}( R \textbackslash{}), and the other is the distance from the point where the pulley is mounted to the vertical where the mass \textbackslash{}( m \textbackslash{}) hangs.

- **Mass \textbackslash{}( m \textbackslash{}):** A blue block labeled \textbackslash{}( m \textbackslash{}) is hanging vertically down from the pulley, indicating it's affected by gravity.

This setup is likely used to study the mechanics of springs, pulleys, and masses in motion, taking into account forces like tension and gravity.

\#\# Dataset metadata

Resonance and Harmonics; Multi-Formula Reasoning, Physical Model Grounding Reasoning

\paragraph{Recorded answer/context.}
D: D: \$22.4\textbackslash{}

\paragraph{Figure/image path.}
/root/proposal\_for\_physic/science-mango/phyx\_mini\_run/phyx\_data/test\_image/235.png

\paragraph{Formalization target.}
create a compiling Lean file with sorry bodies at `PhyXMiniProblems/problem\_phyx\_mini\_0235.lean`. The Lean declarations must preserve the physical quantities, dimensions or dimensional roles, figure labels, governing-law hypotheses, and final relation expressed by this problem.
Use Mathlib/Physlib names found through LeanExplore where available. If a physics API is missing, introduce faithful local abstractions rather than scalar placeholder aliases.

\begin{theorem}[Physics formalization target]
\label{thm:physics:phyx_mini_0235:target}
\lean{PhyXMiniProblems.ProblemPhyXMini0235.problem_phyx_mini_0235}
\uses{def:physics:phyx-mini-0235:phyxminiproblems-problemphyxmini0235-springpulleyoscillator, def:physics:phyx-mini-0235:phyxminiproblems-problemphyxmini0235-matchesproblemreadouts, def:physics:phyx-mini-0235:phyxminiproblems-problemphyxmini0235-matchesscenarioandsuppliedfigure, def:physics:phyx-mini-0235:phyxminiproblems-problemphyxmini0235-hasphysicalparameters, def:physics:phyx-mini-0235:phyxminiproblems-problemphyxmini0235-satisfiesradiusdoubling, def:physics:phyx-mini-0235:phyxminiproblems-problemphyxmini0235-satisfiesstaticequilibrium, def:physics:phyx-mini-0235:phyxminiproblems-problemphyxmini0235-satisfiesnoslipkinematics, def:physics:phyx-mini-0235:phyxminiproblems-problemphyxmini0235-satisfiessoliddiskmomentofinertia, def:physics:phyx-mini-0235:phyxminiproblems-problemphyxmini0235-satisfiessmalloscillationreduction, def:physics:phyx-mini-0235:phyxminiproblems-problemphyxmini0235-roundstodisplayedtenth, def:physics:phyx-mini-0235:phyxminiproblems-problemphyxmini0235-isuniqueclosestdisplayedchoice}
The assigned autoformalize agent should translate this physics problem into Lean declarations in the covered file, with theorem and lemma proof bodies written as `by sorry`.
\end{theorem}
\begin{proof}
This is an autoformalization task, not a proof task. Produce statements that can later be proved by the physics prover without weakening the physical model.
\end{proof}

% --- Archon declaration topology begin ---
\section{Declaration topology}
\begin{definition}[Radius Stage]
  \label{def:physics:phyx-mini-0235:phyxminiproblems-problemphyxmini0235-radiusstage}
  \lean{PhyXMiniProblems.ProblemPhyXMini0235.RadiusStage}
  The pictured pulley before and after its radius is doubled
\end{definition}
\begin{proof}
  This is the declared model definition of Radius Stage.
\end{proof}
\begin{definition}[String Segment]
  \label{def:physics:phyx-mini-0235:phyxminiproblems-problemphyxmini0235-stringsegment}
  \lean{PhyXMiniProblems.ProblemPhyXMini0235.StringSegment}
  The two straight string segments meeting at the pulley
\end{definition}
\begin{proof}
  This is the declared model definition of String Segment.
\end{proof}
\begin{definition}[Orientation]
  \label{def:physics:phyx-mini-0235:phyxminiproblems-problemphyxmini0235-orientation}
  \lean{PhyXMiniProblems.ProblemPhyXMini0235.Orientation}
  The horizontal and vertical orientations visible in the figure
\end{definition}
\begin{proof}
  This is the declared model definition of Orientation.
\end{proof}
\begin{definition}[Figure Object]
  \label{def:physics:phyx-mini-0235:phyxminiproblems-problemphyxmini0235-figureobject}
  \lean{PhyXMiniProblems.ProblemPhyXMini0235.FigureObject}
  Physical objects identifiable in the supplied figure
\end{definition}
\begin{proof}
  This is the declared model definition of Figure Object.
\end{proof}
\begin{definition}[Figure Label]
  \label{def:physics:phyx-mini-0235:phyxminiproblems-problemphyxmini0235-figurelabel}
  \lean{PhyXMiniProblems.ProblemPhyXMini0235.FigureLabel}
  Mathematical labels printed in the supplied figure
\end{definition}
\begin{proof}
  This is the declared model definition of Figure Label.
\end{proof}
\begin{definition}[Mass Idealization]
  \label{def:physics:phyx-mini-0235:phyxminiproblems-problemphyxmini0235-massidealization}
  \lean{PhyXMiniProblems.ProblemPhyXMini0235.MassIdealization}
  Whether a light mechanical component is idealized as massless
\end{definition}
\begin{proof}
  This is the declared model definition of Mass Idealization.
\end{proof}
\begin{definition}[Pulley Shape]
  \label{def:physics:phyx-mini-0235:phyxminiproblems-problemphyxmini0235-pulleyshape}
  \lean{PhyXMiniProblems.ProblemPhyXMini0235.PulleyShape}
  The pulley-shape alternatives relevant to rotational inertia
\end{definition}
\begin{proof}
  This is the declared model definition of Pulley Shape.
\end{proof}
\begin{definition}[Axle Condition]
  \label{def:physics:phyx-mini-0235:phyxminiproblems-problemphyxmini0235-axlecondition}
  \lean{PhyXMiniProblems.ProblemPhyXMini0235.AxleCondition}
  Axle conditions distinguished by the presence of dissipative torque
\end{definition}
\begin{proof}
  This is the declared model definition of Axle Condition.
\end{proof}
\begin{definition}[Pulley Freedom]
  \label{def:physics:phyx-mini-0235:phyxminiproblems-problemphyxmini0235-pulleyfreedom}
  \lean{PhyXMiniProblems.ProblemPhyXMini0235.PulleyFreedom}
  Whether the pulley is free to rotate about its axle
\end{definition}
\begin{proof}
  This is the declared model definition of Pulley Freedom.
\end{proof}
\begin{definition}[String Pulley Contact]
  \label{def:physics:phyx-mini-0235:phyxminiproblems-problemphyxmini0235-stringpulleycontact}
  \lean{PhyXMiniProblems.ProblemPhyXMini0235.StringPulleyContact}
  Contact conditions between the string and pulley rim
\end{definition}
\begin{proof}
  This is the declared model definition of String Pulley Contact.
\end{proof}
\begin{definition}[Oscillation Regime]
  \label{def:physics:phyx-mini-0235:phyxminiproblems-problemphyxmini0235-oscillationregime}
  \lean{PhyXMiniProblems.ProblemPhyXMini0235.OscillationRegime}
  Regimes for the displacement from static equilibrium
\end{definition}
\begin{proof}
  This is the declared model definition of Oscillation Regime.
\end{proof}
\begin{definition}[Release Protocol]
  \label{def:physics:phyx-mini-0235:phyxminiproblems-problemphyxmini0235-releaseprotocol}
  \lean{PhyXMiniProblems.ProblemPhyXMini0235.ReleaseProtocol}
  How the hanging object is prepared for the motion
\end{definition}
\begin{proof}
  This is the declared model definition of Release Protocol.
\end{proof}
\begin{definition}[Spring Pulley Figure]
  \label{def:physics:phyx-mini-0235:phyxminiproblems-problemphyxmini0235-springpulleyfigure}
  \lean{PhyXMiniProblems.ProblemPhyXMini0235.SpringPulleyFigure}
  \uses{def:physics:phyx-mini-0235:phyxminiproblems-problemphyxmini0235-stringsegment, def:physics:phyx-mini-0235:phyxminiproblems-problemphyxmini0235-orientation, def:physics:phyx-mini-0235:phyxminiproblems-problemphyxmini0235-figureobject, def:physics:phyx-mini-0235:phyxminiproblems-problemphyxmini0235-figurelabel}
  Qualitative evidence read directly from the primary image
\end{definition}
\begin{proof}
  This is the declared model definition of Spring Pulley Figure.
\end{proof}
\begin{definition}[Spring Pulley Oscillator]
  \label{def:physics:phyx-mini-0235:phyxminiproblems-problemphyxmini0235-springpulleyoscillator}
  \lean{PhyXMiniProblems.ProblemPhyXMini0235.SpringPulleyOscillator}
  \uses{def:physics:phyx-mini-0235:phyxminiproblems-problemphyxmini0235-radiusstage, def:physics:phyx-mini-0235:phyxminiproblems-problemphyxmini0235-massidealization, def:physics:phyx-mini-0235:phyxminiproblems-problemphyxmini0235-pulleyshape, def:physics:phyx-mini-0235:phyxminiproblems-problemphyxmini0235-axlecondition, def:physics:phyx-mini-0235:phyxminiproblems-problemphyxmini0235-pulleyfreedom, def:physics:phyx-mini-0235:phyxminiproblems-problemphyxmini0235-stringpulleycontact, def:physics:phyx-mini-0235:phyxminiproblems-problemphyxmini0235-oscillationregime, def:physics:phyx-mini-0235:phyxminiproblems-problemphyxmini0235-releaseprotocol, def:physics:phyx-mini-0235:phyxminiproblems-problemphyxmini0235-springpulleyfigure}
  The family of spring--pulley experiments as the unspecified pulley mass M varies. Every real field ending in a unit suffix is a scalar readout in that unit. In particular, no physical primitive is identified with ℝ by an alias. The velocity fields preserve the stated no-slip kinematics, even though the final extremal value depends only on the effective oscillator
\end{definition}
\begin{proof}
  This is the declared model definition of Spring Pulley Oscillator.
\end{proof}
\begin{definition}[Matches Problem Readouts]
  \label{def:physics:phyx-mini-0235:phyxminiproblems-problemphyxmini0235-matchesproblemreadouts}
  \lean{PhyXMiniProblems.ProblemPhyXMini0235.MatchesProblemReadouts}
  \uses{def:physics:phyx-mini-0235:phyxminiproblems-problemphyxmini0235-springpulleyoscillator}
  The numerical readouts stated in the problem, all converted to SI
\end{definition}
\begin{proof}
  This is the declared model definition of Matches Problem Readouts.
\end{proof}
\begin{definition}[Matches Scenario And Supplied Figure]
  \label{def:physics:phyx-mini-0235:phyxminiproblems-problemphyxmini0235-matchesscenarioandsuppliedfigure}
  \lean{PhyXMiniProblems.ProblemPhyXMini0235.MatchesScenarioAndSuppliedFigure}
  \uses{def:physics:phyx-mini-0235:phyxminiproblems-problemphyxmini0235-figureobject, def:physics:phyx-mini-0235:phyxminiproblems-problemphyxmini0235-figurelabel, def:physics:phyx-mini-0235:phyxminiproblems-problemphyxmini0235-springpulleyoscillator}
  The prose assumptions and primary-image incidences. In particular, the image shows that M labels the pulley disk and R labels its radius; it does not show a separate block of mass M
\end{definition}
\begin{proof}
  This is the declared model definition of Matches Scenario And Supplied Figure.
\end{proof}
\begin{definition}[Has Physical Parameters]
  \label{def:physics:phyx-mini-0235:phyxminiproblems-problemphyxmini0235-hasphysicalparameters}
  \lean{PhyXMiniProblems.ProblemPhyXMini0235.HasPhysicalParameters}
  \uses{def:physics:phyx-mini-0235:phyxminiproblems-problemphyxmini0235-radiusstage, def:physics:phyx-mini-0235:phyxminiproblems-problemphyxmini0235-springpulleyoscillator}
  Positivity and nondegeneracy conditions for the physical readouts
\end{definition}
\begin{proof}
  This is the declared model definition of Has Physical Parameters.
\end{proof}
\begin{definition}[Satisfies Radius Doubling]
  \label{def:physics:phyx-mini-0235:phyxminiproblems-problemphyxmini0235-satisfiesradiusdoubling}
  \lean{PhyXMiniProblems.ProblemPhyXMini0235.SatisfiesRadiusDoubling}
  \uses{def:physics:phyx-mini-0235:phyxminiproblems-problemphyxmini0235-springpulleyoscillator}
  The stated modification makes the second pulley radius twice the first
\end{definition}
\begin{proof}
  This is the declared model definition of Satisfies Radius Doubling.
\end{proof}
\begin{definition}[Satisfies Static Equilibrium]
  \label{def:physics:phyx-mini-0235:phyxminiproblems-problemphyxmini0235-satisfiesstaticequilibrium}
  \lean{PhyXMiniProblems.ProblemPhyXMini0235.SatisfiesStaticEquilibrium}
  \uses{def:physics:phyx-mini-0235:phyxminiproblems-problemphyxmini0235-springpulleyoscillator}
  At the undisplaced equilibrium, the spring force balances the hanging object's weight. Gravity therefore fixes the equilibrium extension, though it cancels from the linearized angular frequency
\end{definition}
\begin{proof}
  This is the declared model definition of Satisfies Static Equilibrium.
\end{proof}
\begin{definition}[Satisfies No Slip Kinematics]
  \label{def:physics:phyx-mini-0235:phyxminiproblems-problemphyxmini0235-satisfiesnoslipkinematics}
  \lean{PhyXMiniProblems.ProblemPhyXMini0235.SatisfiesNoSlipKinematics}
  \uses{def:physics:phyx-mini-0235:phyxminiproblems-problemphyxmini0235-radiusstage, def:physics:phyx-mini-0235:phyxminiproblems-problemphyxmini0235-springpulleyoscillator}
  The no-slip rim relation v = R Ω for every modeled trajectory
\end{definition}
\begin{proof}
  This is the declared model definition of Satisfies No Slip Kinematics.
\end{proof}
\begin{definition}[Satisfies Solid Disk Moment Of Inertia]
  \label{def:physics:phyx-mini-0235:phyxminiproblems-problemphyxmini0235-satisfiessoliddiskmomentofinertia}
  \lean{PhyXMiniProblems.ProblemPhyXMini0235.SatisfiesSolidDiskMomentOfInertia}
  \uses{def:physics:phyx-mini-0235:phyxminiproblems-problemphyxmini0235-radiusstage, def:physics:phyx-mini-0235:phyxminiproblems-problemphyxmini0235-springpulleyoscillator}
  The solid-disk moment-of-inertia law I = (1/2) M R²
\end{definition}
\begin{proof}
  This is the declared model definition of Satisfies Solid Disk Moment Of Inertia.
\end{proof}
\begin{definition}[Satisfies Small Oscillation Reduction]
  \label{def:physics:phyx-mini-0235:phyxminiproblems-problemphyxmini0235-satisfiessmalloscillationreduction}
  \lean{PhyXMiniProblems.ProblemPhyXMini0235.SatisfiesSmallOscillationReduction}
  \uses{def:physics:phyx-mini-0235:phyxminiproblems-problemphyxmini0235-radiusstage, def:physics:phyx-mini-0235:phyxminiproblems-problemphyxmini0235-springpulleyoscillator}
  The Newton--Euler reduction for a light, inextensible, non-slipping string on a smooth fixed axle. The pulley contributes the translating effective mass I/R²; the resulting one-dimensional mode has the original spring constant. This is a governing small-oscillation law, not the requested extremal value
\end{definition}
\begin{proof}
  This is the declared model definition of Satisfies Small Oscillation Reduction.
\end{proof}
\begin{lemma}[solid Disk effective Rotational Mass]
  \label{lem:physics:phyx-mini-0235:phyxminiproblems-problemphyxmini0235-soliddisk-effectiverotationalmass}
  \lean{PhyXMiniProblems.ProblemPhyXMini0235.solidDisk_effectiveRotationalMass}
  \uses{def:physics:phyx-mini-0235:phyxminiproblems-problemphyxmini0235-radiusstage, def:physics:phyx-mini-0235:phyxminiproblems-problemphyxmini0235-springpulleyoscillator, def:physics:phyx-mini-0235:phyxminiproblems-problemphyxmini0235-hasphysicalparameters, def:physics:phyx-mini-0235:phyxminiproblems-problemphyxmini0235-satisfiessoliddiskmomentofinertia}
  For a solid disk, I/R² = M/2; hence the radius cancels from the rotational contribution to the effective translating mass
\end{lemma}
\begin{proof}
  The conclusion follows from the governing relations and figure readouts named in the statement of solid Disk effective Rotational Mass.
\end{proof}
\begin{lemma}[angular Frequency squared]
  \label{lem:physics:phyx-mini-0235:phyxminiproblems-problemphyxmini0235-angularfrequency-squared}
  \lean{PhyXMiniProblems.ProblemPhyXMini0235.angularFrequency_squared}
  \uses{def:physics:phyx-mini-0235:phyxminiproblems-problemphyxmini0235-radiusstage, def:physics:phyx-mini-0235:phyxminiproblems-problemphyxmini0235-springpulleyoscillator, def:physics:phyx-mini-0235:phyxminiproblems-problemphyxmini0235-hasphysicalparameters, def:physics:phyx-mini-0235:phyxminiproblems-problemphyxmini0235-satisfiessoliddiskmomentofinertia, def:physics:phyx-mini-0235:phyxminiproblems-problemphyxmini0235-satisfiessmalloscillationreduction}
  The effective PhysLean oscillator obeys ω² = k / (m + M/2) at either radius stage
\end{lemma}
\begin{proof}
  The conclusion follows from the governing relations and figure readouts named in the statement of angular Frequency squared.
\end{proof}
\begin{lemma}[angular Frequency independent of radius Doubling]
  \label{lem:physics:phyx-mini-0235:phyxminiproblems-problemphyxmini0235-angularfrequency-independent-of-radiusdoubling}
  \lean{PhyXMiniProblems.ProblemPhyXMini0235.angularFrequency_independent_of_radiusDoubling}
  \uses{def:physics:phyx-mini-0235:phyxminiproblems-problemphyxmini0235-springpulleyoscillator, def:physics:phyx-mini-0235:phyxminiproblems-problemphyxmini0235-hasphysicalparameters, def:physics:phyx-mini-0235:phyxminiproblems-problemphyxmini0235-satisfiessoliddiskmomentofinertia, def:physics:phyx-mini-0235:phyxminiproblems-problemphyxmini0235-satisfiessmalloscillationreduction}
  Doubling the radius of the same solid disk does not change ω
\end{lemma}
\begin{proof}
  The conclusion follows from the governing relations and figure readouts named in the statement of angular Frequency independent of radius Doubling.
\end{proof}
\begin{definition}[Answer Choice]
  \label{def:physics:phyx-mini-0235:phyxminiproblems-problemphyxmini0235-answerchoice}
  \lean{PhyXMiniProblems.ProblemPhyXMini0235.AnswerChoice}
  The answer labels printed in the source problem
\end{definition}
\begin{proof}
  This is the declared model definition of Answer Choice.
\end{proof}
\begin{definition}[displayed Angular Frequency rad per s]
  \label{def:physics:phyx-mini-0235:phyxminiproblems-problemphyxmini0235-displayedangularfrequency-rad-per-s}
  \lean{PhyXMiniProblems.ProblemPhyXMini0235.displayedAngularFrequency_rad_per_s}
  \uses{def:physics:phyx-mini-0235:phyxminiproblems-problemphyxmini0235-answerchoice}
  The displayed angular-frequency readout for each choice, in radians/s
\end{definition}
\begin{proof}
  This is the declared model definition of displayed Angular Frequency rad per s.
\end{proof}
\begin{definition}[recorded Dataset Answer]
  \label{def:physics:phyx-mini-0235:phyxminiproblems-problemphyxmini0235-recordeddatasetanswer}
  \lean{PhyXMiniProblems.ProblemPhyXMini0235.recordedDatasetAnswer}
  \uses{def:physics:phyx-mini-0235:phyxminiproblems-problemphyxmini0235-answerchoice}
  Dataset metadata recording the supplied answer label
\end{definition}
\begin{proof}
  This is the declared model definition of recorded Dataset Answer.
\end{proof}
\begin{definition}[Rounds To Displayed Tenth]
  \label{def:physics:phyx-mini-0235:phyxminiproblems-problemphyxmini0235-roundstodisplayedtenth}
  \lean{PhyXMiniProblems.ProblemPhyXMini0235.RoundsToDisplayedTenth}
  \uses{def:physics:phyx-mini-0235:phyxminiproblems-problemphyxmini0235-answerchoice, def:physics:phyx-mini-0235:phyxminiproblems-problemphyxmini0235-displayedangularfrequency-rad-per-s}
  Standard rounding to the displayed precision of 0.1 rad/s
\end{definition}
\begin{proof}
  This is the declared model definition of Rounds To Displayed Tenth.
\end{proof}
\begin{definition}[Is Unique Closest Displayed Choice]
  \label{def:physics:phyx-mini-0235:phyxminiproblems-problemphyxmini0235-isuniqueclosestdisplayedchoice}
  \lean{PhyXMiniProblems.ProblemPhyXMini0235.IsUniqueClosestDisplayedChoice}
  \uses{def:physics:phyx-mini-0235:phyxminiproblems-problemphyxmini0235-answerchoice, def:physics:phyx-mini-0235:phyxminiproblems-problemphyxmini0235-displayedangularfrequency-rad-per-s}
  A choice is strictly closer to a value than every other displayed choice
\end{definition}
\begin{proof}
  This is the declared model definition of Is Unique Closest Displayed Choice.
\end{proof}
% --- Archon declaration topology end ---
% --- Archon physics formalization source end ---
